
\newcommand{\bitset}{\{0,1\}}
\newcommand{\eps}{\epsilon}
\newcommand{\light}{\textsf{light}}
\newcommand{\heavy}{\textsf{heavy}}
\renewcommand{\ll}{\textsf{$\ell-\ell$}}
\newcommand{\hh}{\textsf{$h-h$}}
\newcommand{\lh}{\textsf{$\ell-h$}}
\newcommand{\davg}{d_{avg}}
\newcommand{\poly}{\textrm{poly}}

\newcommand{\Ex}{\mathbb{E}}
\newcommand{\Var}{\mathbb{V}}
\newcommand{\eqdef}{\stackrel{\rm def}{=}}



\newcommand{\calA}{{\mathcal{A}}}
\newcommand{\mA}{{\mathcal{A}}}
\newcommand{\calB}{{\mathcal{B}}}
\newcommand{\calC}{{\mathcal{C}}}
\newcommand{\calD}{{\mathcal{D}}}
\newcommand{\calG}{{\mathcal{G}}}
\newcommand{\mG}{{\mathcal{G}}}
\newcommand{\calP}{{\mathcal{P}}}
\newcommand{\calT}{{\mathcal{T}}}

%\newcommand{\trn}{{\widetilde{n}}}
% \newcommand{\trn}{{n_\Delta}}
\newcommand{\trn}{{n'}}
%\newcommand{\trG}{{\widetilde{G}}}
\newcommand{\trG}{{G'}}
%\newcommand{\trm}{{\widetilde{m}}}
\newcommand{\trm}{{m'}}
%\newcommand{\trA}{{\widetilde{\calA}}}
\newcommand{\trA}{{\calA'}}
%\newcommand{\treps}{{\widetilde{\eps}}}
\newcommand{\treps}{{\eps'}}
%\newcommand{\trR}{{\widetilde{R}}}
\newcommand{\trR}{{R'}}
%\newcommand{\tre}{{\widetilde{e}}}
\newcommand{\tre}{e'}


\newcommand{\Sup}{\textrm{Sup}}



\newenvironment{proof}{\smallskip\noindent{\bf Proof:}}%
        {\hspace*{\fill}$\Box$\par}
\newenvironment{proofof}[1]{\smallskip\noindent{\bf Proof of #1:}}%
        {\hspace*{\fill}$\Box$\par}

\newtheorem{theorem}{Theorem}
\newtheorem{corollary}[theorem]{Corollary}
\newtheorem{lemma}{Lemma}
\newtheorem{claim}[lemma]{Claim}
\newtheorem{definition}{Definition}

% I "took" all the \commentts so added D:
\newcommand{\commentt}[1]{ {}}
\newcommand{\talya}[1]{{} }
\newcommand{\reut}[1]{ {}}
\newcommand{\tnew}[1]{{}}

\newcommand{\SodaReb}[1]{{ #1}}

\newcommand{\ICALPCR}[1]{{#1}}


% \newcommand{\ForFuture}[1]{ {\textbf{\color{red} [ForFuture: #1]}}}
\newcommand{\ForFuture}[1]{}

\makeatletter
\def\moverlay{\mathpalette\mov@rlay}
\def\mov@rlay#1#2{\leavevmode\vtop{%
   \baselineskip\z@skip \lineskiplimit-\maxdimen
   \ialign{\hfil$\m@th#1##$\hfil\cr#2\crcr}}}
\newcommand{\charfusion}[3][\mathord]{
    #1{\ifx#1\mathop\vphantom{#2}\fi
        \mathpalette\mov@rlay{#2\cr#3}
      }
    \ifx#1\mathop\expandafter\displaylimits\fi}
\makeatother

\newcommand{\cupdot}{\charfusion[\mathbin]{\cup}{\cdot}}
\newcommand{\bigcupdot}{\charfusion[\mathop]{\bigcup}{\cdot}}

\newenvironment{packed_enum}{
\begin{enumerate}
  \setlength{\itemsep}{0pt}
  \setlength{\parskip}{0pt}
  \setlength{\parsep}{0pt}
}{\end{enumerate}}


 \newcommand{\SelectEdge}{\hyperref[alg:select-edge]{\color{black}{\sf Select-an-Edge}}}
  \newcommand{\TestCfour}{\hyperref[alg:C4]{\color{black}{\sf Test-$C_4$-freeness}}}
 \newcommand{\TestCfive}{\hyperref[alg:C5]{\color{black}{\textsf{Test-$C_5$-freeness}}}}
\newcommand{\TestCsix}{\hyperref[alg:C6]{\color{black}{\textsf{Test-$C_6$-freeness}}}}
\newcommand{\TestCk}{\hyperref[alg:Ck]{\color{black}{\textsf{Test-$C_k$-freeness}}}}
\newcommand{\TestFfree}{\hyperref[alg:Ffree]{\color{black}{\textsf{Test-subgraph-freeness}}}}




 % \textsf{Test-subgraph-freeness}